\documentclass[11pt,a4paper]{article}

% -----------------------------------------------------------------------------
% Required Packages
% -----------------------------------------------------------------------------
\usepackage[utf8]{inputenc}
\usepackage[margin=0.85in]{geometry}
\usepackage{graphicx}
\usepackage{xcolor}
\usepackage{listings}
\usepackage{booktabs}
\usepackage{tabularx}
\usepackage{fancyhdr}
\usepackage{hyperref}
\usepackage{tcolorbox}
\usepackage{caption}
\usepackage{titlesec}

% -----------------------------------------------------------------------------
% Color Palette & Custom Styling
% -----------------------------------------------------------------------------
\definecolor{collegeblue}{RGB}{30, 58, 138}
\definecolor{secondaryblue}{RGB}{30, 41, 59}
\definecolor{accentamber}{RGB}{245, 158, 11}
\definecolor{codebg}{RGB}{248, 250, 252}
\definecolor{codecomment}{RGB}{100, 116, 139}
\definecolor{codekeyword}{RGB}{219, 39, 119}
\definecolor{codestring}{RGB}{16, 185, 129}
\definecolor{lightborder}{RGB}{226, 232, 240}

\hypersetup{
    colorlinks=true,
    linkcolor=collegeblue,
    urlcolor=accentamber,
    citecolor=secondaryblue
}

% Section title styling
\titleformat{\section}
  {\normalfont\Large\bfseries\color{collegeblue}}{\thesection}{1em}{}[\titlerule]
\titleformat{\subsection}
  {\normalfont\large\bfseries\color{secondaryblue}}{\thesubsection}{1em}{}

% Header and Footer setup
\pagestyle{fancy}
\fancyhf{}
\fancyhead[L]{\small \textbf{Web Technology (2304362L)} --- Practical No. 6}
\fancyhead[R]{\small AY 2025--26}
\fancyfoot[L]{\small MIT Academy of Engineering, Pune}
\fancyfoot[R]{\small Page \thepage}
\renewcommand{\headrulewidth}{0.5pt}
\renewcommand{\footrulewidth}{0.5pt}

% -----------------------------------------------------------------------------
% Listings Setup (JavaScript / JSON)
% -----------------------------------------------------------------------------
\lstset{
    backgroundcolor=\color{codebg},
    basicstyle=\ttfamily\footnotesize,
    breaklines=true,
    captionpos=b,
    commentstyle=\color{codecomment}\itshape,
    frame=single,
    rulecolor=\color{lightborder},
    keywordstyle=\color{codekeyword}\bfseries,
    stringstyle=\color{codestring},
    numbers=left,
    numberstyle=\tiny\color{gray},
    numbersep=8pt,
    showstringspaces=false,
    tabsize=2,
    xleftmargin=12pt,
    framexleftmargin=8pt
}

% -----------------------------------------------------------------------------
% Document Body
% -----------------------------------------------------------------------------
\begin{document}

% -----------------------------------------------------------------------------
% TITLE / COVER PAGE
% -----------------------------------------------------------------------------
\begin{titlepage}
    \centering
    \vspace*{0.5cm}
    
    {\Large \textbf{MIT ACADEMY OF ENGINEERING, PUNE}}\\[0.2cm]
    {\large (An Autonomous Institute Affiliated to Savitribai Phule Pune University)}\\[0.2cm]
    {\textbf{School of Computer Engineering}}\\[1.2cm]
    
    \begin{tcolorbox}[colback=collegeblue,colframe=collegeblue,arc=3mm,width=\textwidth,halign=center]
        \vspace{0.3cm}
        {\LARGE \color{white} \textbf{LABORATORY REPORT}}\\[0.3cm]
        {\large \color{accentamber} \textbf{Course: Web Technology (2304362L) | T.Y. B.Tech | AY 2025--26}}
        \vspace{0.3cm}
    \end{tcolorbox}
    
    \vspace{1.2cm}
    
    {\Large \textbf{PRACTICAL NO. 6}}\\[0.4cm]
    {\LARGE \textbf{RESTful CRUD API for Student Records, Deployed and Tested with Postman}}\\[1.5cm]
    
    \begin{minipage}{0.85\textwidth}
        \begin{tabular}{ll}
            \textbf{Student Name:} & Ayush \\
            \textbf{PRN / Roll No.:} & 2304362 \\
            \textbf{Class / Batch:} & T.Y. B.Tech Computer Engineering (Batch B) \\
            \textbf{Course Outcome:} & CO4: Develop and deploy RESTful Web Services \\
            \textbf{Lab Hours:} & 8 Hours \\
            \textbf{Max Marks:} & 10 Marks \\
            \textbf{Date of Performance:} & October 31, 2026 \\
            \textbf{Date of Submission:} & October 31, 2026 \\
        \end{tabular}
    \end{minipage}
    
    \vfill
    
    \begin{tcolorbox}[colback=lightborder!30,colframe=lightborder,arc=2mm,width=0.9\textwidth]
        \centering
        \small
        \textbf{Evaluation Grid (Rubric Alignment)}\\
        \vspace{0.2cm}
        \begin{tabularx}{\textwidth}{XXXX}
            \toprule
            \textbf{Criterion} & \textbf{Max} & \textbf{Marks} & \textbf{Sign} \\
            \midrule
            Correct CRUD endpoints and REST design & 4.0 & & \\
            Status codes, validation and error handling & 2.0 & & \\
            Deployment and Postman collection & 2.0 & & \\
            Endpoint documentation and report & 2.0 & & \\
            \midrule
            \textbf{Total} & \textbf{10.0} & & \\
            \bottomrule
        \end{tabularx}
    \end{tcolorbox}
    
\end{titlepage}

% -----------------------------------------------------------------------------
% 1. AIM & PROBLEM STATEMENT
% -----------------------------------------------------------------------------
\section{Aim and Problem Statement}

\subsection{Aim}
To design, implement, deploy, and rigorously test an enterprise-grade \textbf{RESTful CRUD API for Student Records} using \textbf{Node.js}, \textbf{Express.js}, and \textbf{MongoDB (Mongoose)}, demonstrating semantic HTTP method routing, standardized status codes (200, 201, 400, 404, 500), declarative input validation via \texttt{express-validator}, MongoDB aggregation pipelines, and comprehensive automated testing through an exported \textbf{Postman Collection}.

\subsection{Problem Statement (In Own Words)}
Modern distributed and cross-platform applications require robust, decoupled backend services that expose standardized endpoints for data persistence. Monolithic architectures tightly couple backend computations with HTML presentation logic, limiting multi-client scalability. In this practical assignment, we decouple the presentation layer by architecting a stateless RESTful Web API. The service must execute full Create, Read, Update, and Delete operations against a persistent MongoDB database, enforce database integrity constraints (such as unique PRNs and normalized emails), handle bad client requests gracefully with informative error payloads, and be verifiable via Postman automated assertion scripts.

% -----------------------------------------------------------------------------
% 2. SYSTEM ARCHITECTURE & REST DESIGN
% -----------------------------------------------------------------------------
\section{System Design and Architecture}

The Student Records API is structured around the principles of Representational State Transfer (REST):
\begin{itemize}
    \item \textbf{Stateless Client-Server Communication}: Each HTTP transaction contains all authorization and payload parameters necessary to process the request independently.
    \item \textbf{Resource-Oriented URIs}: Pluralized base noun \texttt{/api/v1/students} identifies the collection, while unique identifiers (\texttt{/api/v1/students/:id}) denote individual resources.
    \item \textbf{Semantic HTTP Verbs}:
    \begin{itemize}
        \item \texttt{POST}: Creates a new student record and emits an HTTP 201 Created status with a \texttt{Location} header.
        \item \texttt{GET}: Idempotent retrieval of the student collection or single student record.
        \item \texttt{PUT}: Updates existing resource state or replaces it.
        \item \texttt{DELETE}: Permanently eliminates the identified resource.
    \end{itemize}
    \item \textbf{Standardized Response Envelope}: All endpoints return consistent JSON objects with \texttt{success}, \texttt{statusCode}, and \texttt{data} or \texttt{message} attributes.
\end{itemize}

% -----------------------------------------------------------------------------
% 3. ENDPOINT SPECIFICATION TABLE
% -----------------------------------------------------------------------------
\section{RESTful Endpoints Specification}

\begin{table}[h!]
\centering
\small
\begin{tabularx}{\textwidth}{llp{7.5cm}l}
\toprule
\textbf{HTTP Method} & \textbf{Endpoint Path} & \textbf{Controller Function / Handler Description} & \textbf{Status Codes} \\
\midrule
\texttt{GET} & \texttt{/health} & Service health check \& MongoDB connection state & 200 \\
\texttt{GET} & \texttt{/api/v1/students} & Retrieve all students with filtering, search, pagination & 200, 500 \\
\texttt{GET} & \texttt{/api/v1/students/:id} & Fetch single student by ObjectId or PRN & 200, 400, 404 \\
\texttt{POST} & \texttt{/api/v1/students} & Validate payload, check uniqueness, insert document & 201, 400 \\
\texttt{PUT} & \texttt{/api/v1/students/:id} & Update existing student fields by ID or PRN & 200, 400, 404 \\
\texttt{DELETE} & \texttt{/api/v1/students/:id} & Remove student record from MongoDB collection & 200, 400, 404 \\
\texttt{GET} & \texttt{/api/v1/students/analytics/stats} & Aggregate metrics: Avg/Max CGPA, Department breakdown & 200 \\
\bottomrule
\end{tabularx}
\caption{Complete RESTful API Endpoint Specification}
\end{table}

% -----------------------------------------------------------------------------
% 4. KEY CODE IMPLEMENTATION
% -----------------------------------------------------------------------------
\section{Key Code Implementation}

\subsection{Mongoose Data Model (\texttt{models/Student.js})}
\begin{lstlisting}[language=JavaScript,caption={Student Schema with validation constraints}]
const mongoose = require('mongoose');

const studentSchema = new mongoose.Schema({
    prn: {
        type: String,
        required: [true, 'PRN is required'],
        unique: true,
        trim: true,
        match: [/^\d{7,8}$/, 'PRN must be 7 to 8 numeric digits']
    },
    name: {
        type: String,
        required: [true, 'Student full name is required'],
        trim: true,
        minlength: [3, 'Name must be at least 3 characters'],
        maxlength: [60, 'Name cannot exceed 60 characters'],
        match: [/^[A-Za-z\s.]+$/, 'Name can only contain alphabets and spaces']
    },
    email: {
        type: String,
        required: [true, 'Email address is required'],
        unique: true,
        lowercase: true,
        match: [/^[a-zA-Z0-9._%+-]+@[a-zA-Z0-9.-]+\.[a-zA-Z]{2,}$/, 'Valid email required']
    },
    department: {
        type: String,
        required: [true, 'Department is required'],
        enum: ['Computer Engineering', 'Information Technology', 'Mechanical Engineering', 
              'Civil Engineering', 'Electronics & Telecommunication', 'Chemical Engineering']
    },
    cgpa: {
        type: Number,
        required: [true, 'CGPA is required'],
        min: [0.0, 'CGPA cannot be less than 0.00'],
        max: [10.0, 'CGPA cannot exceed 10.00']
    }
}, { timestamps: true });

studentSchema.index({ department: 1, year: 1 });
module.exports = { Student: mongoose.model('Student', studentSchema) };
\end{lstlisting}

\subsection{CRUD Controller Logic (\texttt{controllers/studentController.js})}
\begin{lstlisting}[language=JavaScript,caption={Student creation with proactive conflict detection}]
const createStudent = async (req, res, next) => {
    try {
        const { prn, email } = req.body;
        const existing = await Student.findOne({ $or: [{ prn }, { email: email.toLowerCase() }] });
        if (existing) {
            return res.status(400).json({
                success: false,
                statusCode: 400,
                error: 'Conflict / Duplicate Entry',
                message: `Student with this ${existing.prn === prn ? 'PRN' : 'Email'} already exists.`
            });
        }
        const student = await Student.create(req.body);
        res.location(`/api/v1/students/${student._id}`);
        res.status(201).json({
            success: true,
            statusCode: 201,
            message: 'Student record created successfully',
            data: student
        });
    } catch (error) { next(error); }
};
\end{lstlisting}

% -----------------------------------------------------------------------------
% 5. POSTMAN TESTING & RESULTS
% -----------------------------------------------------------------------------
\section{Postman Testing and Output Verification}

The exported Postman collection (\texttt{2304362\_Ayush\_P6\_Postman\_Collection.json}) comprises 12 distinct test cases testing all positive and negative execution paths:

\begin{enumerate}
    \item \textbf{System Health Check (\texttt{GET /health})}: Verifies that the Express process is active and MongoDB connection state is \texttt{Connected} (HTTP 200).
    \item \textbf{Get All Students (\texttt{GET /api/v1/students})}: Retrieves paginated documents array with total count and query pagination metadata (HTTP 200).
    \item \textbf{Create Student (\texttt{POST /api/v1/students})}: Validates JSON payload, persists document, emits HTTP 201, sets \texttt{Location} response header, and populates \texttt{\{\{createdStudentId\}\}} variable.
    \item \textbf{Get Student by Identifier (\texttt{GET /api/v1/students/:id})}: Tests dual-identifier lookup using both PRN (\texttt{2304362}) and MongoDB ObjectId (HTTP 200).
    \item \textbf{Update Student (\texttt{PUT /api/v1/students/:id})}: Updates CGPA to 9.65 and adds technical skills with schema validation (HTTP 200).
    \item \textbf{Delete Student (\texttt{DELETE /api/v1/students/:id})}: Eliminates student record from MongoDB and returns status 200 confirmation.
    \item \textbf{Duplicate PRN Rejection}: Submitting an existing PRN triggers HTTP 400 Bad Request with informative duplicate key message.
    \item \textbf{Validation Failure}: Submitting malformed fields triggers \texttt{express-validator} which rejects the request with HTTP 400 and an array of field errors.
    \item \textbf{Non-Existent Identifier Lookup}: Querying an unknown 24-character hexadecimal ObjectId correctly returns HTTP 404 Not Found.
\end{enumerate}

% -----------------------------------------------------------------------------
% 6. TESTING AND CONCLUSION
% -----------------------------------------------------------------------------
\section{Testing and Conclusion}

\subsection{What Works}
\begin{itemize}
    \item Complete end-to-end CRUD capability against MongoDB using Mongoose models.
    \item Dual-lookup routing supporting both 24-character MongoDB ObjectIds and 7--8 digit student PRNs.
    \item Robust input sanitization and schema enforcement preventing invalid data from entering the database.
    \item Centralized error middleware transforming database CastErrors and E11000 duplicate keys into clean client responses.
    \item Dynamic variable chaining in Postman enabling fully automated end-to-end regression testing.
\end{itemize}

\subsection{Limitations}
\begin{itemize}
    \item Authentication (JWT/Bearer tokens) and role-based permissions were omitted to conform to the basic practical scope.
    \item Rate limiting on API calls could be implemented in future iterations to guard against DoS attacks.
\end{itemize}

\subsection{What I Learned}
\begin{itemize}
    \item Understood the architectural constraints and best practices of RESTful API design.
    \item Mastered MongoDB document schema design, indexing, and aggregation pipelines using Mongoose.
    \item Implemented defensive validation layers using \texttt{express-validator}.
    \item Learned how to configure Postman environment variables, pre-request scripts, and automated test assertions.
\end{itemize}

% -----------------------------------------------------------------------------
% 7. AI USE DECLARATION
% -----------------------------------------------------------------------------
\section{AI Use Declaration}

\begin{table}[h!]
\centering
\small
\begin{tabularx}{\textwidth}{p{4.5cm}X}
\toprule
\textbf{AI Use Declaration} & \textbf{Details} \\
\midrule
\textbf{AI tool and prompt used} & Used Google Antigravity (Gemini 3.8 Flash) with prompt: \textit{"Design and implement a complete RESTful CRUD API for Student Records using ExpressJS, MongoDB, and Mongoose with Postman Collection and validation rules."} \\
\addlinespace
\textbf{One output you checked and how} & Checked the Mongoose duplicate key index error (E11000) and \texttt{express-validator} rule execution by executing automated HTTP test cases with duplicate PRNs and malformed fields via \texttt{test\_api.py}. \\
\addlinespace
\textbf{One thing you corrected or rejected, and why} & Corrected the Mongoose index definition in \texttt{Student.js}: removed duplicate explicit \texttt{studentSchema.index(\{ prn: 1 \})} calls because \texttt{unique: true} on the field already auto-creates the index in Mongoose, eliminating console warnings. \\
\bottomrule
\end{tabularx}
\caption{Mandatory AI Use Declaration (Rubric Aligned)}
\end{table}

\end{document}
